% Source-derived chapter 01: Introduction
% Original source: no-enriques-surface-over-the-integers
\chapter{Introduction}
\label{no-enriques-surface-over-the-integers-chapter-01}
\source{no-enriques-surface-over-the-integers}

\begin{remark}[Source section]
\label{no-enriques-surface-over-the-integers-chapter-01-source}
\source{no-enriques-surface-over-the-integers}
\source{no-enriques-surface-over-the-integers}
This chapter reproduces the corresponding section of the retrieved
LaTeX source, including its definitions, statements, examples, and
proofs. It has not yet been translated into Lean.
\notready
\end{remark}

% The source section title was: Introduction
\mylabel{Introduction}

Which   smooth proper morphism $X\ra\Spec(\ZZ)$ do exist?  This tantalizing question seems to go  back to 
Grothendieck (\cite{Mazur 1986}, page 242), but it is rooted in algebraic number theory:
By   Minkowski's discriminant bound, for each number field $K\neq \QQ$ the corresponding
number ring  is not   \'etale over the ring $R=\ZZ$, so the only examples
in relative dimension $n=0$ are finite sums of $\Spec(\ZZ)$. Moreover, this reduces the general question  to  the
case $\Gamma(X,\O_X)=\ZZ$, such that the structure morphism  $X\ra\Spec(\ZZ)$ is a   contraction of fiber type.

A result of Ogg \cite{Ogg 1966}, which he himself attributes to Tate, asserts that each Weierstra\ss{} equation with integral 
coefficients has discriminant $\Delta\neq \pm 1$. 
Thus there is no family of elliptic curves
$E\ra\Spec(\ZZ)$. In other words, the Deligne--Mumford stack $\shM_{1,1}$ of smooth pointed curves of genus one
has  fiber category $\shM_{1,1}(\ZZ)=\varnothing$.
This was generalized by Fontaine \cite{Fontaine 1985},
who showed that there are no families $A\ra\Spec(\ZZ)$  of non-zero abelian varieties.
In other words, the  fiber category $\shM_{\text{Ab}}(\ZZ)$ for the  stack  
of   abelian varieties contains only  trivial objects.
This was independently established by Abrashkin \cite{Abrashkin 1985}, who also settled the case of abelian surfaces and   threefolds earlier
(\cite{Abrashkin 1976}, \cite{Abrashkin 1977}).

The theory   of relative Picard schemes \cite{FGA V} then shows that there is no family of smooth curves $C\ra\Spec(\ZZ)$
of genus $g\geq 2$. In other words, the Deligne--Mumford stack $\shM_g$ of smooth curves of genus $g$ has fiber category
$\shM_g(\ZZ)=\varnothing$.
Since the ring of integers has trivial Picard and Brauer groups, the only example in relative dimension $n=1$
are finite sums of the projective line $\PP^1_\ZZ$.
Abrashkin \cite{Abrashkin 1990} and Fontaine \cite{Fontaine 1993}   
also established that for each smooth proper $X\ra\Spec(\ZZ)$, the Hodge numbers of the complex fiber $V=X_\CC$
are very restricted, namely
$H^j(V,\Omega^i_{V/\CC}) = 0$ for $i+j\leq 3$ and $i\neq j$.
Note that this ensures that the   stack $\shM_{\text{HK}}$ of hyperk\"ahler varieties has fiber category 
$\shM_{\text{HK}}(\ZZ)=\varnothing$.

The only examples for smooth proper schemes over the integers that easily come to mind
are  the projective space $\PP^n$ and schemes   stemming from this,
for example sums, products,   blowing ups with respect to linear centers, or the Hilbert scheme 
$X=\Hilb^n_{\PP^2/\ZZ}$.
Similar examples arise from $G/P$, where $G$ is a Chevalley group  and $P$ is a  parabolic subgroup,
and from   torus embeddings $\operatorname{Temb}_\Delta$, stemming from a  fan of regular cones $\sigma\in\Delta$.
I am not aware of any other construction. It would be interesting to understand the case of   elliptic surfaces over $\PP^1$.

In light of the Enriques classification of algebraic surfaces, it is natural to consider  families
$Y\ra \Spec(\ZZ)$ of smooth proper surfaces with $c_1=0$. 
These are the abelian surfaces, bielliptic surfaces, K3 surfaces and Enriques surfaces.
They can be distinguished by their second Betti number, which takes the   values $b_2=6,2,22,10$.
The former two are ruled out by the Hodge numbers $h^{0,1}\geq 1$ of the complex fiber. Similarly, K3 surfaces do not 
occur, by $h^{2,0}=1$.
So only the Enriques surfaces remain, which over the complex numbers indeed have $h^{i,j}=0$ for $i+j\leq 3$ and $i\neq j$.

At first glance, the family of  K3 coverings $X\ra Y$ seems to show that there is no   family of Enriques surfaces.
There is, however, a  crucial loophole: At prime $p=2$, the K3 cover can degenerate into a \emph{K3-like covering}, and then
the fiber $X\otimes\FF_2$ becomes singular. In light of this, it actually seems    possible that
such families of Enriques surfaces exist. The main result of this paper asserts that it does not happen:

\begin{maintheorem}
\source{no-enriques-surface-over-the-integers}
\notready
{\rm (see Thm.\ \ref{no enriques over integers})}
The   stack $\shM_\Enr$ of Enriques surfaces
has fiber category $\shM_\Enr(\ZZ)=\varnothing$. In other words, 
there is no smooth proper family  of Enriques surfaces over the ring $R=\ZZ$. 
\end{maintheorem}

Note that the finiteness of the fiber categories $\shM_\Enr(R)$ for rings of integers $R$ in   number fields  was recently
established by Takamatsu \cite{Takamatsu 2019}. A similar finiteness holds for abelian varieties
by Faltings \cite{Faltings 1983} and Zarhin \cite{Zarhin 1985}, and also for K3 surfaces by Andr\'e \cite{Andre 1996}.
Results of this type are often refereed to as the \emph{Shafarevich Conjecture} \cite{Shafarevich 1963}

The proof of our non-existence result is long and indirect, and is given in the final Section \ref{Proof}.
It is actually a  consequence of the following:

\begin{maintheorem}
\source{no-enriques-surface-over-the-integers}
\notready
{\rm (see Thm.\ \ref{no non-exceptional})}
There is no Enriques surface $Y$ over   $k=\FF_2$  that is non-exceptional and has
constant Picard scheme $\Pic_{Y/k}=(\ZZ^{\oplus 10}\oplus\ZZ/2\ZZ)_k$.
\end{maintheorem}

The  \emph{exceptional Enriques surfaces} where introduced by Ekedahl and Shepherd-Barron \cite{Ekedahl; Shepherd-Barron 2004}.
They indeed can be discarded from our considerations, as I recently showed in \cite{Schroeer 2021b} that they do not admit a lifting
to the ring $W_2(\FF_2)=\ZZ/4\ZZ$.
The above result will be deduced from an explicit classification of \emph{geometrically rational elliptic surfaces} $\phi:J\ra\PP^1$
over $k=\FF_2$ that have constant Picard scheme, and satisfy 
certain additional technical conditions stemming from the theory of Enriques surfaces:

\begin{maintheorem}
\source{no-enriques-surface-over-the-integers}
\notready
{\rm (see Thm.\ \ref{classification non-zero j} and Thm.\ \ref{classification zero j})}
Up to isomorphisms, there are exactly eleven Weierstra\ss{}
equations $y^2+a_1xy + \ldots =x^3+a_2x^2+\ldots$ with coefficients $a_i\in \FF_2[t]$
that define a geometrically rational elliptic surface $\phi:J\ra\PP^1$  
with constant Picard scheme $\Pic_{J/\FF_2}$  having at most one rational point $a\in \PP^1$ where $J_a$ is semistable or supersingular.
\end{maintheorem}

For our Enriques surfaces $Y$ over the prime field $k=\FF_2$ with constant Picard scheme,
this narrows down the possible configurations of fibers over the three rational points $t=0,1,\infty$
in genus-one fibrations $\varphi:Y\ra\PP^1$, by passing to the jacobian fibration $\phi:J\ra\PP^1$ and using 
a result of Liu, Lorenzini and Raynaud \cite{Liu; Lorenzini; Raynaud 2005}.

It turns out that there are fifteen possible configurations, 
where the  additional four cases come from quasielliptic fibrations. In all but one situation there is precisely one ``non-reduced''
Kodaira symbol, from the list $\II^*,\III^*,\IV^*,\I^*_4, \I_2^*, \I_1^*$.
We then make an extensive combinatorial analysis for the possible interaction of pairs of genus-one fibrations
$\varphi,\psi:Y\ra\PP^1$ whose intersection number takes the minimal value  $\varphi^\inv(\infty)\cdot\psi^\inv(\infty)=4$,
in the spirit of Cossec and Dolgachev \cite{Cossec; Dolgachev 1989}, Chapter III, \S 5.
This only becomes feasible after  discarding the   exceptional Enriques surfaces,
which roughly speaking have ``too few'' fibrations and ``too large'' degenerate fibers.

The paper is organized as follows:
In Section \ref{Local system} we treat the group of numerical classes of invertible
sheaves as a local system $\Num_{X/k}$, that is, a sheaf on the \'etale site of the ground field $k$. 
Here the main result is that if it is constant,
then every projective contraction of the base-change to $k^\alg$ descends to a contraction of
$X$.
Geometric consequences for smooth surfaces $X=S$ are examined in more detail in Section \ref{Contractions}.
Then we analyze curves of canonical type in Section \ref{Curves}. They 
occur in genus-one fibrations, and we stress   arithmetical aspects.  
In Section \ref{Enriques surfaces} we review the theory of Enriques surfaces over algebraically closed fields, including
the notion of exceptional Enriques surfaces. Section \ref{Families} contains a detailed study of families  of Enriques surfaces,
in particular over the ring of integers.
In Section \ref{Geometrically rational} 
we collect basic facts on geometrically rational elliptic surfaces, again from an arithmetical point of view.
In Section \ref{Counting points}   we turn to counting points over finite fields with the Weil Conjectures,
and give some formulas that apply to Enriques surfaces and geometrically rational elliptic surfaces
whose local system $\Num_{X/k}$ is constant.
Section \ref{Multiple fibers} contains an observation on the behavior of reduced fibers in passing from
an elliptic fibration to its jacobian fibration.
Section \ref{Passage to jacobian} contains a main technical result: If $Y$ is an   Enriques surface over $k=\FF_2$
with constant Picard scheme, then the resulting geometrically rational surfaces $J$, arising as jacobian for
elliptic fibrations, also has constant Picard scheme.
Using this, we classify in Section \ref{Classification jacobian} all Weierstra\ss{} equations with coefficients from $\FF_2[t]$ that
possibly could describe such $J$. In Section \ref{Possible configurations}, we make a list of the possible fiber configurations
over the rational points $t=0,1,\infty$ that could arise for the Enriques surface $Y$,
where we also take into account quasielliptic fibrations.
This list is narrowed down to two configurations in Sections  \ref{Elimination I4*} and \ref{Elimination III*}.
In Section \ref{Restricted configurations} we   work again over arbitrary algebraically closed ground fields $k$, and show that
Enriques surfaces all whose genus-one fibrations are subject to certain severe combinatorial restrictions do not exist.
In the final Section \ref{Proof} we combine all these observations and deduce  our main results.

\begin{acknowledgement}
The research was supported by the Deutsche Forschungsgemeinschaft by the grant  \emph{SCHR 671/6-1 Enriques-Mannigfaltigkeiten}.
It was also conducted  in the framework of the   research training group
\emph{GRK 2240: Algebro-geometric Methods in Algebra, Arithmetic and Topology}.
I wish to thank Will Sawin for useful remarks, and Matthias Sch\"utt for helpful comments, in particular for 
pointing out  \cite{Shioda 1992} and \cite{Schuett; Shioda 2019},
which correct two entries in the classification of Mordell--Weil lattices for rational elliptic surfaces.
I would also like to thank the referees for very  careful and thorough reading and many valuable suggestions, which helped to improve the paper
and remove mistakes. In particular, for pointing out a defective argument in the first version concerning   two-sections in
what are  now Propositions \ref{no I4*+R} and   \ref{no I4*+R second}.
\end{acknowledgement}


%===========================================================
